\subsection*{2.2. 消灭圈}

\subsubsection*{2.2.1.}
设 $S$ 为 Hensel 特质, $s,\eta,\bar s,\bar\eta$ 的意义如 XIII, 0.2.5. 置 $\Lambda=\mathbb Z/k$, 其中 $k$ 与 $S$ 的剩余特征互素. 设 $\mathbb A^{n+1}_S$ 为 $S$ 上的标准 $(n+1)$ 维仿射空间, $X$ 为 $\mathbb A^{n+1}_S$ 中由二次方程 $Q=0$ 定义的闭子概形, 且该方程模一个一致化元非零:
\[
Q(\underline x)=\sum_{i\leq j}a_{ij}x_ix_j+\sum_i b_ix_i+c.
\]
将 $\mathbb A^{n+1}_S$ 视为 $\mathbb P^{n+1}_S$ 中无穷远超平面 $H$ 的补集. 则 $X$ 是从由齐次方程
\begin{equation}\tag{2.2.1.1}
\sum_{i\leq j}a_{ij}x_ix_j+\sum_i b_ix_iz+cz^2=0
\end{equation}
定义的二次超曲面 $X_1\subset\mathbb P^{n+1}_S$ 中去掉 $Y=X_1\cap H$ 得到的.

作以下假设:

\begin{enumerate}
\item[(a)] $Y$ 为 $S$ 上的光滑二次超曲面, 即二次型 $\sum_{i\leq j}a_{ij}x_ix_j$ 为普通二次型;

\item[(b)] 二次超曲面 $X_s$ 有奇点, 即 $X_{\bar s}$ 为二次锥.
\end{enumerate}
\footnote{印本在 (b) 中写为 $X'_s$. 等价条件以及紧随其后的奇点定义表明这里应为 $X_s$.}

以 $x_0$ 表示 $X_s$ 的奇点, 以 $f:X\to S$ 表示投影.

\subsubsection*{注 2.2.2.}
若 $x_0$ 为有理点, 则经过平移可将其置于坐标原点. 若 $x_0=0$, 则 $b_i$ 和 $c$ 均可被一致化元整除.

$X_s$ 中由方程 $\partial Q/\partial x_i$ 定义的子概形 $A$ 集中于 $x_0$. 除了 $k(s)$ 的特征为 $2$ 且 $n+1$ 为偶数这一例外情形以外, 此概形在 $k(s)$ 上的次数为一, 从而 $x_0$ 为有理点. 在例外情形中, 它在 $k(s)$ 上分歧, 秩为 $2$, 故 $k(x_0)$ 或为 $k(s)$, 或为 $k(s)$ 的纯不可分二次扩张. 关于 $A$ 的这些断言, 转到 $k(s)$ 的代数闭包 $k(\bar s)$ 后最易检验.

\subsubsection*{命题 2.2.3.}
\underline{在 2.2.1 的假设下, 态射}
\[
H^i(X_{\bar\eta},\Lambda)
 \xrightarrow{\ \sim\ }
\mathbb H^i\bigl(X_{\bar s},R\psi_{\bar\eta}(\Lambda)\bigr)
 \xrightarrow{\ \sim\ }
R^i\psi_{\bar\eta}(\Lambda)_{x_0}
\]
\[
H_c^i(X_{\bar\eta},\Lambda)
 \xleftarrow{\ \sim\ }
\mathbb H_c^i\bigl(X_{\bar s},R\psi_{\bar\eta}(\Lambda)\bigr)
 \xleftarrow{\ \sim\ }
\mathbb H^i_{\{x_0\}}\bigl(X_{\bar s},R\psi_{\bar\eta}(\Lambda)\bigr)
\]
\underline{为同构.}

左边箭头为同构由 2.1.8.6 及 2.1.10.5 得出; 关于右边箭头的断言是 2.1.2 的推论. 事实上, 有区分三角
\[
(\Lambda\text{ 在 }X_{\bar s})[0]
 \longrightarrow R\psi_{\bar\eta}(\Lambda)
 \longrightarrow R\Phi_{\bar\eta}(\Lambda),
\]
而 2.1.2 适用于 $(\Lambda\text{ 在 }X_{\bar s})$, 同时 $R\Phi_{\bar\eta}(\Lambda)$ 的支集为 $x_0$.

\subsubsection*{推论 2.2.4.}
\underline{若几何一般纤维 $X_{\bar\eta}$ 有奇点, 即它仍是二次锥, 则所有层 $R^i\Phi(\Lambda)$ 均为零.}

由 2.1.2,
\[
\Lambda=H^0(X_{\bar s},\Lambda)
 \xrightarrow{\ \sim\ }H^0(X_{\bar\eta},\Lambda),
\qquad
0=H^i(X_{\bar s},\Lambda)=H^i(X_{\bar\eta},\Lambda)
\quad(i>0).
\]
再应用长正合列 XIII, 2.1.8.9(由 XIII, 2.1.9 可用).
\footnote{印本两处 $H^o$ 中均使用字母 $o$; 上同调次数应为 $0$.}